% Appendix. Sources: refs/docs/ts_benchmark.md §3 (synthesis rules);
% HTML appendix A (results by scenario, tables/scenarios.tex).

\section{Synthesis of the AmphionTS-Bench synthetic part}
\label{app:synthesis}

\paragraph{Levels.}
Each utterance is first normalized to $-25$\,dBFS and then scaled according to the signal-to-interference ratio of the condition, with an additional $\pm2$\,dB per interferer; the final mixture is peak-normalized to 0.5.

\paragraph{Timing.}
The target speaks only complete sentences.
The last sentence of an interferer may be cut off at the end of the clip, with a 30\,ms fade-out, to approximate segments cut by voice activity detection.

\paragraph{Noise floors and rooms.}
Noise floors come from WHAM!~\cite{wichern2019wham} (real ambient noise such as cafés, restaurants and streets), stationary AudioSet and Freesound excerpts distributed with the DNS Challenge~\cite{dubey2023icassp}, and white and pink noise, at $-65$ to $-35$\,dBFS.
Room responses come from measured responses in OpenSLR 28, simulated small, medium and large rooms, and the simulated responses of the DNS Challenge.

\paragraph{Background sounds.}
Background conditions use MUSAN~\cite{snyder2015musan} noise and music, AudioSet traffic, FSD50K~\cite{fonseca2022fsd50k} environmental events and crowds, synthetic babble from 5--20 test-split speakers, and Mandarin pop songs with accompaniment from MIR-1K, used only for evaluation.
% TODO: add citations for MIR-1K, AudioSet and OpenSLR 28 (not yet in refs INDEX).

\paragraph{Diagnostic conditions.}
Twelve subsets outside the sections cover 0\,dB enrollment contamination, clipping, 3.4\,kHz narrowband audio, $\mu$-law coding, and stacked outdoor and office conditions.

\section{Results by scenario}
\label{app:scenarios}

\Cref{tab:scenarios} lists the complete results on \tsbench grouped by scenario.
Each synthetic group changes a single condition on the same base scenes.

{\small
% Generated from the published HTML report (refs/internal/html-report/, table 20;
% data: data.json["v12"]["groups"]). Bold marks the best value in each row.
\begin{longtable}{@{}>{\raggedright\arraybackslash}p{5.4cm}crccc@{}}
  \caption{TS-ASR results on \tsbench by scenario (\%). Target-present rows give CER / empty output; target-absent rows give the false-alarm rate. Bold marks the best value in each row.}
  \label{tab:scenarios}\\
  \toprule
  Scenario & Target & Samples & AmphionASR 2.0 & CocktailASR-1 & CoT \\
  \midrule
  \endfirsthead
  \toprule
  Scenario & Target & Samples & AmphionASR 2.0 & CocktailASR-1 & CoT \\
  \midrule
  \endhead
  \bottomrule
  \endfoot
  \multicolumn{6}{@{}l}{\textit{Degraded enrollment}} \\
  Target present (degraded enrollment) & present & 1,200 & \best{13.3 / 1.9} & 33.1 / 1.8 & 51.7 / 6.0 \\
  Target absent (degraded enrollment) & absent & 1,200 & \best{5.8} & 76.3 & 50.2 \\
  Only 2--4 s of enrollment recorded (target present) & present & 450 & \best{13.5 / 1.8} & 26.9 / 0.2 & 31.8 / 1.1 \\
  Only 2--4 s of enrollment recorded (target absent) & absent & 450 & \best{4.2} & 81.1 & 63.8 \\
  \midrule
  \multicolumn{6}{@{}l}{\textit{Entry, exit and number of interferers}} \\
  Target present (entry, exit, number of interferers) & present & 1,920 & \best{15.1 / 3.0} & 46.8 / 4.3 & 68.5 / 7.5 \\
  Target absent (others speaking) & absent & 960 & \best{4.9} & 78.4 & 73.0 \\
  Target finishes, another speaker follows (no overlap) & present & 360 & \best{3.8 / 3.3} & 10.7 / 0.0 & 38.6 / 1.1 \\
  \midrule
  \multicolumn{6}{@{}l}{\textit{Short segments}} \\
  Target absent $\cdot$ 1.2--6 s segment (others speaking) & absent & 480 & \best{3.3} & 11.0 & 37.5 \\
  Target absent $\cdot$ another speaker says one short sentence & absent & 600 & \best{13.7} & 37.8 & 69.8 \\
  Target alone, one short sentence (control) & present & 420 & \best{4.2 / 0.5} & 4.8 / 0.9 & 4.4 / 0.2 \\
  Enrolled speaker says one short sentence (quiet, far-field, noisy) & present & 2,400 & \best{6.5 / 2.4} & 14.3 / 6.7 & 12.2 / 0.4 \\
  \midrule
  \multicolumn{6}{@{}l}{\textit{No one speaking}} \\
  No speech & absent & 720 & \best{0.1} & 10.8 & 58.1 \\
  Noise only, no target & absent & 360 & \best{2.2} & 14.7 & 54.7 \\
  Singing only (songs with vocals) & absent & 120 & \best{0.0} & 4.2 & 14.2 \\
  \midrule
  \multicolumn{6}{@{}l}{\textit{Environmental noise and recording chain}} \\
  Target present (environment and device) & present & 960 & \best{13.5 / 1.8} & 26.1 / 0.4 & 32.0 / 1.6 \\
  Target absent (environment and device) & absent & 720 & \best{3.1} & 83.1 & 63.5 \\
  Stress panel $\cdot$ noise-floor type and level (target present) & present & 1,200 & \best{9.9 / 2.6} & 19.9 / 0.2 & 24.6 / 1.9 \\
  Stress panel $\cdot$ noise-floor type and level (target absent) & absent & 1,200 & \best{2.1} & 79.2 & 61.2 \\
  Stress panel $\cdot$ level and recording chain (target present) & present & 360 & \best{13.3 / 2.8} & 26.4 / 1.1 & 30.8 / 1.9 \\
  Stress panel $\cdot$ level and recording chain (target absent) & absent & 360 & \best{1.4} & 77.5 & 59.7 \\
  Background noise (target present) & present & 840 & \best{10.7 / 2.6} & 26.3 / 0.6 & 26.5 / 2.5 \\
  Background noise (target absent) & absent & 840 & \best{2.0} & 77.6 & 59.0 \\
  \midrule
  \multicolumn{6}{@{}l}{\textit{Far field and room reverberation}} \\
  Stress panel $\cdot$ room response (target present) & present & 600 & \best{13.7 / 3.2} & 29.4 / 1.3 & 33.5 / 3.2 \\
  Stress panel $\cdot$ room response (target absent) & absent & 600 & \best{0.8} & 74.8 & 56.2 \\
  Separate room responses for target and interferers (target present) & present & 600 & \best{14.7 / 3.0} & 31.1 / 2.0 & 35.0 / 4.2 \\
  Separate room responses for target and interferers (target absent) & absent & 600 & \best{2.0} & 71.8 & 55.3 \\
  Near/far-field mismatch between enrollment and mixture (target present) & present & 480 & \best{14.2 / 3.5} & 32.2 / 1.7 & 41.2 / 3.1 \\
  Near/far-field mismatch between enrollment and mixture (target absent) & absent & 480 & \best{3.3} & 74.0 & 57.3 \\
  \midrule
  \multicolumn{6}{@{}l}{\textit{Real meetings, phone calls and web conversations}} \\
  AliMeeting far-field meetings (target present) & present & 115 & \best{6.8 / 0.9} & 38.2 / 0.0 & 36.9 / 0.0 \\
  AliMeeting far-field meetings (target absent) & absent & 100 & \best{1.0} & 100.0 & 97.0 \\
  AISHELL-4 far-field meetings (target present) & present & 141 & \best{15.0 / 5.0} & 27.7 / 0.0 & 23.0 / 0.0 \\
  AISHELL-4 far-field meetings (target absent) & absent & 140 & \best{5.0} & 100.0 & 97.1 \\
  MagicData-RAMC phone calls (target present) & present & 100 & \best{13.3 / 0.0} & 16.1 / 0.0 & 16.0 / 0.0 \\
  MagicData-RAMC phone calls (target absent) & absent & 100 & \best{5.0} & 96.0 & 83.0 \\
  Emilia web conversations, pseudo-labels (target present) & present & 148 & \best{10.2 / 2.0} & 19.7 / 0.0 & 16.6 / 0.0 \\
  Emilia web conversations, pseudo-labels (target absent) & absent & 150 & \best{2.7} & 74.0 & 50.0 \\
  AliMeeting near/far-field pairs (target present) & present & 460 & \best{7.1 / 2.2} & 36.3 / 0.0 & 31.5 / 0.2 \\
  AliMeeting near/far-field pairs (target absent) & absent & 400 & \best{2.0} & 99.8 & 90.5 \\
  Real meetings $\cdot$ enrollment contains another voice (target present) & present & 62 & \best{15.6 / 1.6} & 60.8 / 0.0 & 64.1 / 0.0 \\
  Real meetings $\cdot$ enrollment contains another voice (target absent) & absent & 77 & \best{1.3} & 97.4 & 96.1 \\
  AliMeeting near/far-field pairs $\cdot$ enrollment contains another voice (target present) & present & 164 & \best{18.5 / 4.9} & 94.1 / 0.0 & 87.0 / 0.6 \\
  AliMeeting near/far-field pairs $\cdot$ enrollment contains another voice (target absent) & absent & 224 & \best{2.2} & 100.0 & 96.9 \\
\end{longtable}

}
